% Источники: exam/materials/преподаватель/2020/23-03-2020-MODEL L6(2).doc; exam/materials/преподаватель/2020/27-04-2020-Model_L10_3_.pdf; exam/materials/моделирование.pdf, разделы 29 и 34; exam/materials/прошлогодние ответы/Modelirovanie_1_-3.pdf, вопрос 29; GitHub HumsterProgrammer/iu9-notes/8sem/Моделирование/Моделирование_лекция-13_26-03-26.md.
\section{Анализ функциональных зависимостей между наблюдаемыми одновременно характеристиками объекта. Простая линейная регрессия и простой корреляционный анализ.}

Когда мы говорим о функциональной зависимости между переменными, мы подразумеваем, что значение одной величины (зависимой переменной) может быть предсказано или выражено через значения других величин (независимых переменных).

Под \textbf{регрессией} понимают условное математическое ожидание случайной переменной $y \in Y$ при условии, что другая переменная $x \in X$ приняла конкретное значение, например $x=A$:
\[
  E[Y \mid X=A].
\]

Тогда моделью линейной регрессии является модель, в которой теоретическое среднее значение наблюдаемой величины $y$ является линейной комбинацией независимых переменных.

\textbf{Простая линейная регрессия} используется, когда есть одна независимая переменная, и зависимость между $Y$ и $X$ предполагается линейной:
\[
  Y = a + bX + \varepsilon,
\]
где $Y$ — зависимая переменная; $X$ — независимая переменная; $a$ — свободный член; $b$ — коэффициент регрессии; $\varepsilon$ — случайная (неустранимая) ошибка, показывает, что в реальных задачах мы не можем идеально определить зависимость.

\textbf{Множественная регрессия} применяется, когда имеется несколько независимых переменных, влияющих на целевую переменную:
\[
  Y = \beta_0 + \beta_1 X_1 + \cdots + \beta_n X_n + \varepsilon.
\]
Можно показать, что
\[
  \hat{\beta} = (X^T X)^{-1}X^TY,
\]
где $X$ — размерности (кол-во наблюдений, кол-во признаков), элементы — значения признаков на данном наблюдении.

\subsection*{Корреляционный анализ}

\textbf{Корреляция}:
\begin{enumerate}
  \item мера степени и направления связи между значениями двух переменных;
  \item статистический показатель вероятности связи между двумя переменными, измеренными в количественной шкале.
\end{enumerate}

\textbf{Корреляционный анализ} — проверка гипотез о связях между переменными с использованием коэффициентов корреляции.

Направление связи определяется прямым или обратным соотношением значений двух переменных: если возрастанию значений одной переменной соответствует возрастание значений другой переменной, то взаимосвязь называется прямой (положительной); если возрастанию значений одной переменной соответствует убывание значений другой переменной, то взаимосвязь является обратной (отрицательной).

Показателем направления связи является знак коэффициента корреляции.

\textbf{Коэффициент корреляции ($r$)}:
\begin{enumerate}
  \item количественная мера силы и направления вероятностной взаимосвязи двух переменных; принимает значения в диапазоне от $-1$ до $+1$;
  \item мера прямой или обратной пропорциональности между двумя переменными;
  \item двумерная описательная статистика, количественная мера взаимосвязи (совместимой изменчивости) двух переменных.
\end{enumerate}

Коэффициенты: Пирсона, Спирмана, Кенделла. Для определения тесноты связи двух качественных признаков, каждый из которых состоит только из двух групп (альтернативные признаки), используют коэффициенты ассоциации, контингенции Пирсона ($\varphi$-коэффициент), коллигации.

Для переменных с интервальной и номинальной шкалой используется коэффициент корреляции Пирсона. Если, по меньшей мере, одна из двух переменных имеет порядковую шкалу, либо не является нормально распределённой, используется ранговая корреляция Спирмана. Применение коэффициента Кендалла предпочтительно, если в исходных данных имеются выбросы.

\subsection*{Пример}

Пары точек с координатами $(x,y)$ наносятся на координатную сетку. По облаку точек получают предварительное представление о рассеивании: возможна функциональная зависимость, стохастическая зависимость с выраженным трендом или отсутствие явно выраженного тренда из-за ошибок измерения и изменения результата. Поэтому задача разбивается на две подзадачи: установить вид тренда (линейный или нелинейный) и подтвердить тесноту связи.

\clearpage
